% Generated draft: requires mathematical and source review.
\documentclass[11pt]{article}
\usepackage[margin=25mm]{geometry}
\usepackage{fontspec}
\setmainfont{Latin Modern Roman}
\usepackage{amsmath,amssymb,mathtools,unicode-math}
\setmathfont{Latin Modern Math}
\usepackage{enumitem,needspace,xurl,hyperref}
\setlength{\parindent}{0pt}
\setlength{\parskip}{5pt}
\setlength{\emergencystretch}{3em}
\newcommand{\N}{\mathbb{N}}
\newcommand{\Z}{\mathbb{Z}}
\newcommand{\Q}{\mathbb{Q}}
\newcommand{\R}{\mathbb{R}}
\newcommand{\C}{\mathbb{C}}
\newcommand{\F}{\mathbb{F}}
\newcommand{\A}{\mathbb{A}}
\newcommand{\PP}{\mathbb{P}}
\newcommand{\E}{\mathbb{E}}
\newcommand{\eps}{\varepsilon}
\newcommand{\dd}{\,\mathrm{d}}
\begin{document}
\section*{UnsolvedMath ID: 800200}
\textbf{Source problem code:} catalogue entry
\subsection{Definitions and mathematical statement}
A projective K3 surface is a smooth, simply connected complex projective surface with trivial canonical bundle. A polarization on such a surface \(S\) is an ample line bundle \(H\); its genus is conventionally defined by
\[
g=1+\frac{H^2}{2}.
\]
For an integer \(r\geq 2\), write \(S^{[r]}=\operatorname{Hilb}^r(S)\) for the Hilbert scheme of zero-dimensional subschemes of \(S\) of length \(r\). A polarization on \(S^{[r]}\) means an ample line bundle \(L\) on that Hilbert scheme. A polarized deformation of \((S^{[r]},L)\) is a deformation of the variety together with a line bundle extending \(L\), ample on the fibers under consideration. Here a family is called locally complete for polarized deformations if, at a general base point, its classifying map covers a neighborhood in the local deformation space of the polarized fiber.

\textbf{Problem (Beauville, 2010, Question 9).} For the genera listed in the source as \(g=1,2,3,\ldots\), describe the general projective deformation of \(S^{[2]}\), where \(S\) is a polarized K3 surface of genus \(g\), allowing a suitable choice of polarization \(L\) on \(S^{[2]}\). As a weaker alternative, find further examples of locally complete projective families of such deformations. Ask the same questions for \(S^{[r]}\) for each integer \(r\geq 3\).

\textbf{Unresolved qualifications.} With the conventional definition above, ampleness gives \(H^2>0\), so the source's genus-\(1\) case cannot be interpreted as an ample polarization. It would require another convention, potentially involving a nef square-zero class or a genus-one fibration, but the supplied record does not specify one. Thus the conventional formulation applies to \(g\geq 2\), while the intended genus-\(1\) formulation remains unresolved. The source also leaves the polarization \(L\), its numerical type, and the meaning of ``general'' unspecified. The polarized-deformation interpretation and the definition of local completeness above make explicit a natural framework, but the record does not determine whether this is exactly the intended deformation space. It does not specify whether \(S\) itself is general in its polarized moduli space or whether its polarization must be primitive.
\subsection{Short English statement}
Beauville asks for descriptions of general projective deformations of Hilbert schemes of points on polarized K3 surfaces, or at least further projective families that locally realize all relevant deformations. The question begins with two points and extends to every \(r\geq 3\). The genus-\(1\) case and the precise polarization and deformation conventions require clarification.
\subsection{Sources}
\begin{itemize}
\item Local UnsolvedMath registry, supplied complete source record: permanent ID \(800200\), exact source problem code \texttt{catalogue entry}, title ``Holomorphic Symplectic Geometry --- Question 9.'' Attribution supplied by the registry: Beauville, 2010. The registry labels the problem \emph{open}; this is a supplied status, not an independently verified mathematical or literature assessment.
\item Beauville, \emph{Holomorphic Symplectic Geometry: A Problem List} (2010), Question 9. Supplied source URL: \url{https://arxiv.org/abs/1002.4321}. The supplied provenance records access on 2026-10-16 and extraction by \texttt{arxiv\_latex}; neither access nor extraction has been independently verified for this draft.
\item Supplied upstream revision: \texttt{d371ee917ac9647fbb06cd6fc529deec6f3dda82}. The record was created and updated on 2026-10-16. No subsequent literature assessment is supplied.
\end{itemize}
\end{document}
